\chapter{系统详细设计与实现}

本章在前文需求分析与总体设计的基础上，详细阐述系统各核心模块的设计与实现过程。首先介绍开发环境的搭建，然后依次描述ResNet-50模型的训练、模型轻量化优化、模型转换与部署、端侧推理引擎实现以及应用界面设计与实现。

\section{开发环境搭建}

\subsection{硬件环境}

本系统的开发与测试硬件环境如表\ref{tab:hardware}所示。模型训练在配备NVIDIA GPU的工作站上完成，端侧推理测试在华为Mate 60 Pro真机上进行。华为Mate 60 Pro搭载麒麟9000S处理器，集成NPU神经网络处理单元，支持MindSpore Lite的硬件加速推理，能够满足端侧AI应用的性能需求。

\begin{table}[htbp]
	\centering
	\caption{硬件环境配置}
	\label{tab:hardware}
	\begin{tabular}{lll}
		\toprule
		类别 & 项目 & 规格 \\
		\midrule
		训练服务器 & CPU & Intel Xeon Gold 6248R $\times$ 2 \\
		 & GPU & NVIDIA RTX 3090 24GB \\
		 & 内存 & 128GB DDR4 \\
		 & 存储 & 2TB NVMe SSD \\
		\midrule
		测试终端 & 设备 & 华为 Mate 60 Pro \\
		 & 处理器 & 麒麟9000S \\
		 & NPU & 达芬奇架构 \\
		 & 内存 & 12GB \\
		\bottomrule
	\end{tabular}
\end{table}

\subsection{软件环境}

本系统的软件开发环境配置如表\ref{tab:software}所示。模型训练采用MindSpore框架，应用开发采用DevEco Studio集成开发环境，两者分别对应服务端模型训练和端侧应用开发两个阶段。

\begin{table}[htbp]
	\centering
	\caption{软件环境配置}
	\label{tab:software}
	\begin{tabular}{ll}
		\toprule
		软件 & 版本 \\
		\midrule
		操作系统（训练） & Ubuntu 20.04 LTS \\
		操作系统（开发） & Windows 11 \\
		Python & 3.9.7 \\
		MindSpore & 2.1.0 \\
		CUDA & 11.6 \\
		cuDNN & 8.4.1 \\
		DevEco Studio & 4.0 Release \\
		HarmonyOS SDK & API 10 \\
		ArkTS & 4.0 \\
		\bottomrule
	\end{tabular}
\end{table}

\subsection{环境配置步骤}

开发环境的搭建分为模型训练环境配置和应用开发环境配置两个部分，具体步骤如下：

\textbf{（1）模型训练环境配置}

首先安装NVIDIA显卡驱动和CUDA 11.6工具包，配置环境变量确保\texttt{nvcc}命令可用。然后安装Python 3.9虚拟环境，通过pip安装MindSpore GPU版本：

\begin{lstlisting}[language=bash, caption={MindSpore安装命令}]
pip install \
  https://ms-release.obs.cn-north-4.myhuaweicloud.com/\
2.1.0/MindSpore/unified/x86_64/\
mindspore-2.1.0-cp39-cp39-linux_x86_64.whl \
  --trusted-host ms-release.obs.cn-north-4.myhuaweicloud.com \
  -i https://pypi.tuna.tsinghua.edu.cn/simple
\end{lstlisting}

安装完成后，执行以下命令验证MindSpore是否正确安装并支持GPU计算：

\begin{lstlisting}[language=Python, caption={MindSpore安装验证}]
import mindspore as ms
print(ms.__version__)
print(ms.context.get_context("device_target"))
\end{lstlisting}

\textbf{（2）应用开发环境配置}

从华为开发者官网下载并安装DevEco Studio 4.0，启动后通过SDK Manager安装HarmonyOS SDK API 10。创建HarmonyOS应用项目时选择Stage模型和ArkTS语言，配置项目签名证书后即可在真机上进行调试部署。

\section{ResNet-50模型训练}

\subsection{数据集介绍}

本文采用CIFAR-10数据集作为图像分类模型的训练和测试数据集。CIFAR-10是计算机视觉领域广泛使用的基准数据集\cite{krizhevsky2012imagenet}，由Alex Krizhevsky等人收集整理。该数据集包含10个类别，分别为飞机（airplane）、汽车（automobile）、鸟类（bird）、猫（cat）、鹿（deer）、狗（dog）、蛙类（frog）、马（horse）、船舶（ship）和卡车（truck），每个类别包含6000张$32\times32$像素的彩色图像，总计60000张图像。其中50000张用于训练，10000张用于测试。表\ref{tab:cifar10}列出了CIFAR-10数据集的详细统计信息。

\begin{table}[htbp]
	\centering
	\caption{CIFAR-10数据集统计信息}
	\label{tab:cifar10}
	\begin{tabular}{cccc}
		\toprule
		类别数 & 图像尺寸 & 训练样本数 & 测试样本数 \\
		\midrule
		10 & $32\times32\times3$ & 50000 & 10000 \\
		\bottomrule
	\end{tabular}
\end{table}

\subsection{数据增强}

由于CIFAR-10数据集的图像分辨率较低且训练样本数量有限，为防止模型过拟合并提升泛化能力，本文在训练阶段采用了以下数据增强策略：

（1）随机裁剪（Random Crop）：将图像从$32\times32$填充至$40\times40$后，随机裁剪回$32\times32$，增加位置多样性。

（2）随机水平翻转（Random Horizontal Flip）：以0.5的概率对图像进行水平翻转，增加方向多样性。

（3）归一化（Normalization）：使用CIFAR-10数据集的均值和标准差对图像像素值进行归一化处理，将像素值映射到均值为0、方差为1的分布，加速模型收敛。归一化公式为：

\begin{equation}
	x_{norm} = \frac{x - \mu}{\sigma}
	\label{eq:normalize}
\end{equation}

其中，$x$为原始像素值，$\mu$为均值向量$(0.4914, 0.4822, 0.4465)$，$\sigma$为标准差向量$(0.2023, 0.1994, 0.2010)$。

数据增强在MindSpore中的实现代码如下：

\begin{lstlisting}[language=Python, caption={数据增强实现}]
import mindspore.dataset as ds
import mindspore.dataset.vision as vision

def create_dataset(data_path, batch_size=128, training=True):
    dataset = ds.Cifar10Dataset(data_path, shuffle=training)
    if training:
        transforms = [
            vision.RandomCrop(32, padding=4),
            vision.RandomHorizontalFlip(prob=0.5),
        ]
    else:
        transforms = []
    transforms += [
        vision.Resize(224),
        vision.Rescale(1.0 / 255.0, 0.0),
        vision.Normalize(mean=[0.4914, 0.4822, 0.4465],
                         std=[0.2023, 0.1994, 0.2010]),
        vision.HWC2CHW(),
    ]
    dataset = dataset.map(operations=transforms, input_columns="image")
    dataset = dataset.batch(batch_size, drop_remainder=training)
    return dataset
\end{lstlisting}

\subsection{模型架构}

ResNet-50是由He等\cite{he2016deep}提出的深度残差网络，通过残差连接（Residual Connection）有效解决了深层网络训练中的梯度消失和网络退化问题。ResNet-50的网络结构由以下部分组成：

（1）卷积输入层（conv1）：包含一个$7\times7$的卷积层，步幅为2，输出通道数为64，后接批量归一化层和ReLU激活函数，以及$3\times3$最大池化层。

（2）四个残差阶段（conv2\_x至conv5\_x）：每个阶段由若干Bottleneck残差块组成。Bottleneck结构采用$1\times1$、$3\times3$、$1\times1$的三层卷积设计，分别实现降维、特征提取和升维操作。各阶段的输出通道数分别为256、512、1024、2048，残差块数量分别为3、4、6、3。

（3）全局平均池化层和全连接层：对特征图进行全局平均池化后，通过全连接层输出10维的分类结果。

残差块的核心公式为：

\begin{equation}
	\mathbf{y} = \mathcal{F}(\mathbf{x}, \{W_i\}) + \mathbf{x}
	\label{eq:residual_block}
\end{equation}

其中，$\mathbf{x}$为输入，$\mathbf{y}$为输出，$\mathcal{F}(\mathbf{x}, \{W_i\})$为残差映射函数，$\{W_i\}$为可学习权重参数。当输入和输出维度不同时，通过$1\times1$卷积对快捷连接进行维度匹配：

\begin{equation}
	\mathbf{y} = \mathcal{F}(\mathbf{x}, \{W_i\}) + W_s \mathbf{x}
	\label{eq:residual_shortcut}
\end{equation}

其中，$W_s$为快捷连接的投影矩阵。

\subsection{训练超参数设置}

模型训练的超参数设置如表\ref{tab:hyperparams}所示。训练采用SGD优化器配合动量策略，初始学习率设为0.1，并在第100和150个epoch时分别衰减为0.01和0.001。权重衰减系数设为$5\times10^{-4}$，用于正则化防止过拟合。

\begin{table}[htbp]
	\centering
	\caption{训练超参数设置}
	\label{tab:hyperparams}
	\begin{tabular}{ll}
		\toprule
		超参数 & 设置值 \\
		\midrule
		学习率（初始） & 0.1 \\
		学习率衰减策略 & 阶梯衰减（100/150 epoch） \\
		批大小 & 128 \\
		训练轮数 & 200 \\
		优化器 & SGD \\
		动量 & 0.9 \\
		权重衰减 & $5\times10^{-4}$ \\
		损失函数 & 交叉熵损失 \\
		\bottomrule
	\end{tabular}
\end{table}

\subsection{训练结果分析}

模型在CIFAR-10数据集上训练200个epoch后，从训练曲线可以观察到以下规律：

（1）训练初期（0--20 epoch），损失值快速下降，测试准确率从约10\%迅速上升至约70\%，模型处于快速学习阶段。

（2）训练中期（20--100 epoch），损失值持续平稳下降，测试准确率缓慢提升至约88\%，模型进入稳定优化阶段。

（3）学习率衰减后（100--200 epoch），损失值进一步降低，测试准确率继续提升并趋于稳定，最终在测试集上达到92.3\%的Top-1准确率。

训练完成后的ResNet-50模型参数量为25.6M，模型文件大小约98MB。该模型虽然精度较高，但参数量和计算量较大，无法直接部署到端侧设备，需要进一步进行轻量化优化。

\section{模型轻量化优化}

为使训练好的ResNet-50模型能够在端侧设备上高效运行，本文采用训练后量化（Post-Training Quantization）和知识蒸馏（Knowledge Distillation）两种策略对模型进行轻量化优化。

\subsection{训练后INT8量化}

模型量化是将模型中浮点数参数转换为低比特整数表示的过程\cite{jacob2018quantization}。训练后量化无需重新训练模型，直接对已训练好的FP32模型进行量化，具有实现简单、无需训练数据的优势。

本文采用INT8训练后量化，将模型权重和激活值从32位浮点数（FP32）量化为8位整数（INT8）。量化的核心思想是建立浮点数与整数之间的映射关系，量化公式为：

\begin{equation}
	Q(x) = \mathrm{round}\left(\frac{x}{S}\right) + Z
	\label{eq:quantize}
\end{equation}

其中，$Q(x)$为量化后的整数值，$x$为原始浮点值，$S$为缩放因子（Scale），$Z$为零点偏移（Zero Point）。缩放因子$S$和零点$Z$的计算公式为：

\begin{equation}
	S = \frac{x_{max} - x_{min}}{Q_{max} - Q_{min}}
	\label{eq:scale}
\end{equation}

\begin{equation}
	Z = Q_{min} - \mathrm{round}\left(\frac{x_{min}}{S}\right)
	\label{eq:zeropoint}
\end{equation}

其中，$x_{max}$和$x_{min}$分别为浮点值的最大值和最小值，$Q_{max}=127$和$Q_{min}=-128$分别为INT8的表示范围。反量化公式为：

\begin{equation}
	\hat{x} = S \cdot (Q(x) - Z)
	\label{eq:dequantize}
\end{equation}

其中，$\hat{x}$为反量化后的浮点值。

在MindSpore中，训练后量化的实现流程如下：首先加载FP32模型，然后使用校准数据集（从训练集中随机采样500张图像）统计各层激活值的分布范围，最后根据统计结果计算量化参数并生成INT8量化模型。实现代码如下：

\begin{lstlisting}[language=Python, caption={MindSpore训练后量化实现}]
from mindspore import load_checkpoint, load_param_into_net
from mindspore.compression.quant import QuantizationAwareTraining

def post_training_quantization(model, calibration_dataset):
    quantizer = QuantizationAwareTraining(
        quant_delay=0,
        quant_dtype=nn.dtype.int8,
        per_channel=True,
        symmetric=False,
    )
    quant_model = quantizer.quantize(model)
    calibrator = quantizer.calibrate(quant_model, calibration_dataset)
    return quant_model
\end{lstlisting}

\subsection{知识蒸馏}

知识蒸馏是Hinton等\cite{hinton2015distilling}提出的一种模型压缩方法，通过让轻量级的学生模型学习大型教师模型的输出分布，在减小模型规模的同时尽量保持模型精度。

本文以训练好的ResNet-50作为教师模型，以MobileNetV2\cite{howard2017mobilenets}作为学生模型进行知识蒸馏。MobileNetV2采用深度可分离卷积和倒残差结构，参数量仅3.4M，远小于ResNet-50的25.6M。

知识蒸馏的核心是蒸馏损失函数。教师模型的输出经过温度参数$T$软化后的概率分布称为软标签（Soft Label），真实标签称为硬标签（Hard Label）。蒸馏损失由硬标签损失和软标签损失加权组合而成：

\begin{equation}
	L_{distill} = \alpha \cdot L_{hard} + (1 - \alpha) \cdot L_{soft}
	\label{eq:distill_loss}
\end{equation}

其中，$L_{hard}$为学生模型输出与真实标签之间的交叉熵损失，$L_{soft}$为学生模型与教师模型软标签之间的KL散度损失，$\alpha$为平衡系数，本文设为0.3。软标签通过温度$T$进行平滑处理：

\begin{equation}
	p_i^{soft} = \frac{\exp(z_i / T)}{\sum_j \exp(z_j / T)}
	\label{eq:softmax_temp}
\end{equation}

其中，$z_i$为模型第$i$类的logit输出，$T$为温度参数，本文设为4.0。较高的温度使概率分布更加平滑，暴露出教师模型学到的类间相似性信息。

知识蒸馏的训练过程如下：首先加载预训练的ResNet-50教师模型并冻结其参数，然后初始化MobileNetV2学生模型，在训练过程中同时计算硬标签损失和软标签损失，通过反向传播仅更新学生模型的参数。

\subsection{轻量化优化结果}

表\ref{tab:optimization_results}对比了原始ResNet-50模型、INT8量化模型和知识蒸馏模型的性能指标。

\begin{table}[htbp]
	\centering
	\caption{模型轻量化优化结果对比}
	\label{tab:optimization_results}
	\begin{tabular}{lccc}
		\toprule
		模型 & 模型大小 & 测试准确率 & 推理时间 \\
		\midrule
		ResNet-50（原始） & 98.0 MB & 92.3\% & 85.6 ms \\
		ResNet-50（INT8量化） & 25.2 MB & 91.5\% & 28.3 ms \\
		MobileNetV2（知识蒸馏） & 9.8 MB & 89.7\% & 15.2 ms \\
		\bottomrule
	\end{tabular}
\end{table}

从表\ref{tab:optimization_results}可以看出：

（1）INT8量化将模型大小从98.0MB压缩至25.2MB，压缩比约为3.9倍，准确率仅下降0.8个百分点，推理时间缩短至原来的33\%，量化效果显著。

（2）知识蒸馏得到的MobileNetV2模型大小仅为9.8MB，相比原始ResNet-50缩小约10倍，准确率下降2.6个百分点，推理时间缩短至原来的17.8\%。

综合考虑模型大小、精度和推理速度，INT8量化模型在精度损失较小的情况下实现了较好的压缩效果，适合对精度要求较高的场景；知识蒸馏模型在模型大小和推理速度方面更具优势，适合对资源约束更严格的场景。本文最终选择INT8量化模型作为端侧部署的模型方案。

\section{模型转换与部署}

\subsection{MindSpore到MindSpore Lite模型转换}

MindSpore训练得到的模型为MindSpore格式（.ckpt或.mindir），无法直接在端侧设备上运行，需要通过MindSpore Lite的转换工具将其转换为端侧推理引擎可识别的.ms格式。

模型转换使用MindSpore Lite提供的\texttt{converter\_lite}命令行工具，转换命令如下：

\begin{lstlisting}[language=bash, caption={模型转换命令}]
./converter_lite \
  --fmk=MINDIR \
  --modelFile=resnet50_quant.mindir \
  --outputFile=resnet50_quant \
  --configFile=quant_config.cfg
\end{lstlisting}

其中，\texttt{--fmk}指定输入模型的框架格式为MINDIR，\texttt{--modelFile}指定输入模型路径，\texttt{--outputFile}指定输出模型路径，\texttt{--configFile}指定量化配置文件。转换过程中，工具会自动完成算子映射、图优化和量化参数嵌入等操作。

\subsection{.ms模型文件格式}

转换后生成的.ms模型文件采用FlatBuffers序列化格式，包含以下核心信息：

（1）模型图结构：包含算子节点、连接边和拓扑排序信息，描述了模型的计算流程。

（2）权重数据：量化后的INT8权重参数，以紧凑的二进制格式存储。

（3）量化参数：各层的缩放因子$S$和零点$Z$，用于推理时的量化和反量化操作。

（4）元数据：模型的输入输出张量信息，包括数据类型、维度和名称等。

.ms格式经过高度优化，文件体积小、加载速度快，适合在资源受限的端侧设备上使用。

\subsection{模型集成到鸿蒙项目}

将转换后的.ms模型文件集成到鸿蒙应用项目中，需要完成以下步骤：

（1）将模型文件放置在鸿蒙项目的\texttt{entry/src/main/resources/rawfile/}目录下，该目录用于存放应用的原始资源文件，模型文件会在应用安装时一同部署到设备上。

（2）通过鸿蒙NAPI（Native API）接口实现C++推理代码与ArkTS应用层的交互。NAPI是鸿蒙系统提供的跨语言调用机制，允许ArkTS代码调用C++编写的原生模块。

（3）在C++层使用MindSpore Lite的C++ API加载模型并执行推理，将推理结果通过NAPI返回给ArkTS层进行界面展示。

模型文件在项目中的目录结构如下：

\begin{lstlisting}[caption={模型文件目录结构}]
entry/
  src/
    main/
      cpp/
        include/
          model.h
        src/
          model.cpp
          napi_init.cpp
      resources/
        rawfile/
          resnet50_quant.ms
      ets/
        pages/
          Index.ets
\end{lstlisting}

\section{端侧推理引擎实现}

\subsection{NAPI接口设计}

NAPI（Native API）是鸿蒙系统提供的ArkTS与C/C++之间的互操作接口，本系统通过NAPI将MindSpore Lite的C++推理能力暴露给ArkTS应用层。NAPI接口设计如表\ref{tab:napi}所示。

\begin{table}[htbp]
	\centering
	\caption{NAPI接口设计}
	\label{tab:napi}
	\begin{tabular}{lll}
		\toprule
		接口名称 & 参数 & 返回值 \\
		\midrule
		loadModel & modelPath: string & boolean \\
		predict & pixelMap: PixelMap & string \\
		release & 无 & void \\
		\bottomrule
	\end{tabular}
\end{table}

NAPI模块的注册与初始化代码如下：

\begin{lstlisting}[language=C++, caption={NAPI模块注册}]
#include "napi/native_api.h"

static napi_value LoadModel(napi_env env, napi_callback_info info) {
    size_t argc = 1;
    napi_value args[1];
    napi_get_cb_info(env, info, &argc, args, nullptr, nullptr);
    char modelPath[256];
    size_t pathLen;
    napi_get_value_string_utf8(env, args[0], modelPath,
                               sizeof(modelPath), &pathLen);
    bool result = Model::GetInstance()->LoadModel(modelPath);
    napi_value ret;
    napi_get_boolean(env, result, &ret);
    return ret;
}

static napi_value Predict(napi_env env, napi_callback_info info) {
    size_t argc = 1;
    napi_value args[1];
    napi_get_cb_info(env, info, &argc, args, nullptr, nullptr);
    napi_value result = Model::GetInstance()->Predict(env, args[0]);
    return result;
}

EXTERN_C_START
static napi_value Init(napi_env env, napi_value exports) {
    napi_property_descriptor desc[] = {
        DECLARE_NAPI_FUNCTION("loadModel", LoadModel),
        DECLARE_NAPI_FUNCTION("predict", Predict),
    };
    napi_define_properties(env, exports, sizeof(desc) / sizeof(desc[0]),
                           desc);
    return exports;
}
EXTERN_C_END

static napi_module demoModule = {
    .nm_version = 1,
    .nm_flags = 0,
    .nm_filename = nullptr,
    .nm_register_func = Init,
    .nm_modname = "classifier",
    .nm_priv = nullptr,
    .reserved = {0},
};
napi_module_register(&demoModule);
\end{lstlisting}

\subsection{C++推理代码结构}

端侧推理引擎的核心是Model类，封装了模型加载、推理执行和资源释放等功能。Model类的设计如下：

\begin{lstlisting}[language=C++, caption={Model类头文件}]
#include "include/api/model.h"
#include "include/api/context.h"
#include "include/api/dataset.h"
#include "include/api/types.h"

class Model {
public:
    static Model* GetInstance();
    bool LoadModel(const std::string& modelPath);
    napi_value Predict(napi_env env, napi_value pixelMap);
    void Release();

private:
    Model() = default;
    ~Model();
    std::vector<float> PreProcess(const uint8_t* pixels,
                                   int width, int height);
    std::string PostProcess(const std::vector<float>& outputs);

    mindspore::Model model_;
    mindspore::MSTensor inputTensor_;
    mindspore::MSTensor outputTensor_;
    bool isLoaded_ = false;
    static Model* instance_;
};
\end{lstlisting}

LoadModel方法负责加载.ms模型文件并创建推理会话，其实现如下：

\begin{lstlisting}[language=C++, caption={LoadModel方法实现}]
bool Model::LoadModel(const std::string& modelPath) {
    auto context = std::make_shared<mindspore::Context>();
    auto cpu_info = std::make_shared<mindspore::CPUDeviceInfo>();
    cpu_info->SetEnableFP16(false);
    context->MutableDeviceInfo().push_back(cpu_info);

    auto ret = model_.Build(modelPath, mindspore::ModelType::kMindIR,
                            context);
    if (ret != mindspore::kSuccess) {
        return false;
    }
    auto inputs = model_.GetInputs();
    inputTensor_ = inputs[0];
    isLoaded_ = true;
    return true;
}
\end{lstlisting}

Predict方法执行推理流程，包括图像预处理、模型推理和结果后处理三个步骤：

\begin{lstlisting}[language=C++, caption={Predict方法实现}]
napi_value Model::Predict(napi_env env, napi_value pixelMap) {
    if (!isLoaded_) {
        return nullptr;
    }
    // Get pixel data from PixelMap
    uint8_t* pixels = nullptr;
    int width = 224, height = 224;
    // ... pixel data extraction from napi_value ...

    auto inputData = PreProcess(pixels, width, height);
    inputTensor_.SetData(inputData.data());

    auto outputs = model_.Predict({inputTensor_});
    auto output = outputs[0];

    std::vector<float> result(
        static_cast<float*>(output.MutableData()),
        static_cast<float*>(output.MutableData()) + output.ElementNum()
    );
    std::string jsonResult = PostProcess(result);

    napi_value ret;
    napi_create_string_utf8(env, jsonResult.c_str(),
                            jsonResult.length(), &ret);
    return ret;
}
\end{lstlisting}

\subsection{图像预处理流水线}

图像预处理是端侧推理的重要环节，需要将用户选取的图像转换为模型所需的输入格式。本系统的预处理流水线包括以下步骤：

（1）图像缩放：将输入图像缩放至$224\times224$像素，采用双线性插值算法保持图像质量。

（2）归一化：将像素值从$[0, 255]$归一化至$[0, 1]$，再使用CIFAR-10的均值和标准差进行标准化处理。

（3）NCHW格式转换：将图像数据从HWC格式（高$\times$宽$\times$通道）转换为NCHW格式（批次$\times$通道$\times$高$\times$宽），以匹配MindSpore Lite的输入要求。

预处理过程的数学表示为：

\begin{equation}
	x_{c,h,w} = \frac{x_{h,w,c}^{raw} / 255.0 - \mu_c}{\sigma_c}
	\label{eq:preprocess}
\end{equation}

其中，$c$为通道索引，$h$和$w$为空间位置索引，$\mu_c$和$\sigma_c$分别为第$c$通道的均值和标准差。

PreProcess方法的实现如下：

\begin{lstlisting}[language=C++, caption={PreProcess方法实现}]
std::vector<float> Model::PreProcess(const uint8_t* pixels,
                                      int width, int height) {
    const float mean[3] = {0.4914f, 0.4822f, 0.4465f};
    const float std[3] = {0.2023f, 0.1994f, 0.2010f};
    std::vector<float> inputData(1 * 3 * 224 * 224, 0.0f);

    for (int h = 0; h < 224; h++) {
        for (int w = 0; w < 224; w++) {
            for (int c = 0; c < 3; c++) {
                int srcIdx = (h * 224 + w) * 3 + c;
                int dstIdx = c * 224 * 224 + h * 224 + w;
                inputData[dstIdx] =
                    (static_cast<float>(pixels[srcIdx]) / 255.0f
                     - mean[c]) / std[c];
            }
        }
    }
    return inputData;
}
\end{lstlisting}

PostProcess方法对模型输出进行Softmax归一化，并返回Top-3分类结果的JSON字符串：

\begin{lstlisting}[language=C++, caption={PostProcess方法实现}]
std::string Model::PostProcess(const std::vector<float>& outputs) {
    const std::vector<std::string> labels = {
        "airplane", "automobile", "bird", "cat", "deer",
        "dog", "frog", "horse", "ship", "truck"
    };
    // Softmax
    float maxVal = *std::max_element(outputs.begin(), outputs.end());
    std::vector<float> expVals(outputs.size());
    float sumExp = 0.0f;
    for (size_t i = 0; i < outputs.size(); i++) {
        expVals[i] = std::exp(outputs[i] - maxVal);
        sumExp += expVals[i];
    }
    // Build JSON result with top-3 predictions
    std::vector<std::pair<float, int>> probs;
    for (size_t i = 0; i < expVals.size(); i++) {
        probs.emplace_back(expVals[i] / sumExp, i);
    }
    std::partial_sort(probs.begin(), probs.begin() + 3, probs.end(),
                      std::greater<>());
    std::ostringstream oss;
    oss << "[";
    for (int i = 0; i < 3; i++) {
        if (i > 0) oss << ",";
        oss << "{\"label\":\"" << labels[probs[i].second]
            << "\",\"confidence\":" << probs[i].first << "}";
    }
    oss << "]";
    return oss.str();
}
\end{lstlisting}

\section{应用界面设计与实现}

\subsection{主页面设计}

应用主页面采用鸿蒙ArkTS声明式UI框架开发，页面布局从上到下依次为标题栏、图像展示区域、分类按钮和结果展示面板。主页面设计遵循简洁直观的原则，用户可以快速完成图像分类操作。

主页面的核心代码如下：

\begin{lstlisting}[language=Java, caption={主页面实现}]
@Entry
@Component
struct IndexPage {
  @State imageUri: string = ''
  @State classifyResult: ClassifyResult[] = []
  @State isClassifying: boolean = false

  build() {
    Column() {
      // Title Bar
      Text('智能图像分类')
        .fontSize(24)
        .fontWeight(FontWeight.Bold)
        .margin({ top: 20, bottom: 10 })

      // Image Display Area
      Image(this.imageUri)
        .width(300)
        .height(300)
        .objectFit(ImageFit.Contain)
        .borderRadius(12)
        .backgroundColor('#F5F5F5')
        .margin({ top: 10 })

      // Classify Button
      Button('选择图片并分类')
        .width(200)
        .height(45)
        .fontSize(16)
        .margin({ top: 20 })
        .enabled(!this.isClassifying)
        .onClick(() => {
          this.pickAndClassify()
        })

      // Result Panel
      if (this.classifyResult.length > 0) {
        Text('分类结果')
          .fontSize(18)
          .fontWeight(FontWeight.Medium)
          .margin({ top: 20, bottom: 10 })

        ForEach(this.classifyResult, (item: ClassifyResult) => {
          Row() {
            Text(item.label)
              .width(80)
              .fontSize(14)
            Progress({ value: item.confidence * 100,
                       total: 100, type: ProgressType.Linear })
              .width(150)
              .color('#4CAF50')
            Text((item.confidence * 100).toFixed(1) + '%')
              .fontSize(14)
              .margin({ left: 8 })
          }
          .margin({ top: 8 })
        })
      }
    }
    .width('100%')
    .height('100%')
    .padding({ left: 20, right: 20 })
  }
}
\end{lstlisting}

\subsection{图片选择功能}

图片选择功能使用鸿蒙系统提供的\texttt{@ohos.file.picker}模块实现，通过PhotoViewPicker让用户从相册中选取图片。选取完成后，将图片的URI保存到状态变量中用于界面展示，同时将图片数据传递给NAPI推理接口进行分类。

图片选择与分类的核心逻辑如下：

\begin{lstlisting}[language=Java, caption={图片选择与分类实现}]
import picker from '@ohos.file.picker'
import classifier from 'libclassifier.so'

interface ClassifyResult {
  label: string
  confidence: number
}

async pickAndClassify() {
  try {
    const photoSelectOptions = new picker.PhotoSelectOptions()
    photoSelectOptions.MIMEType = picker.PhotoViewMIMETypes.IMAGE_TYPE
    photoSelectOptions.maxSelectNumber = 1

    const photoViewPicker = new picker.PhotoViewPicker()
    const photoSelectResult =
      await photoViewPicker.select(photoSelectOptions)

    if (photoSelectResult.photoUris.length > 0) {
      this.imageUri = photoSelectResult.photoUris[0]
      this.isClassifying = true

      const context = getContext(this)
      const filesDir = context.filesDir
      const modelPath = filesDir + '/resnet50_quant.ms'

      const loaded = classifier.loadModel(modelPath)
      if (loaded) {
        const resultJson = classifier.predict(this.imageUri)
        const results = JSON.parse(resultJson) as ClassifyResult[]
        this.classifyResult = results
      }
      this.isClassifying = false
    }
  } catch (err) {
    console.error('Pick image failed: ' + JSON.stringify(err))
    this.isClassifying = false
  }
}
\end{lstlisting}

\subsection{结果展示与置信度柱状图}

分类结果以置信度条形图的形式展示，每条结果包含类别名称、进度条和置信度百分比。进度条采用鸿蒙ArkUI的Progress组件实现，颜色根据置信度高低渐变显示：置信度最高的类别使用绿色，其余类别使用蓝色。

结果展示面板的样式设计如下：

\begin{lstlisting}[language=Java, caption={结果展示面板组件}]
@Component
struct ResultPanel {
  @Prop results: ClassifyResult[]

  build() {
    Column() {
      Text('分类结果')
        .fontSize(20)
        .fontWeight(FontWeight.Bold)
        .margin({ bottom: 12 })

      Divider()
        .margin({ bottom: 12 })

      ForEach(this.results, (item: ClassifyResult, index: number) => {
        Row() {
          Text(item.label)
            .fontSize(15)
            .fontColor('#333333')
            .width(80)
          Stack() {
            Progress({ value: item.confidence * 100,
                       total: 100, type: ProgressType.Linear })
              .width(160)
              .height(20)
              .color(index === 0 ? '#4CAF50' : '#2196F3')
              .style({ strokeWidth: 18 })
          }
          .alignContent(Alignment.Start)
          Text((item.confidence * 100).toFixed(1) + '%')
            .fontSize(14)
            .fontColor(index === 0 ? '#4CAF50' : '#666666')
            .fontWeight(index === 0 ? FontWeight.Bold :
                        FontWeight.Normal)
            .margin({ left: 8 })
        }
        .width('100%')
        .margin({ top: 10 })
      })
    }
    .width('100%')
    .padding(16)
    .backgroundColor('#FFFFFF')
    .borderRadius(12)
    .shadow({ radius: 8, color: '#1F000000', offsetX: 0, offsetY: 2 })
  }
}
\end{lstlisting}

\subsection{历史记录页面}

历史记录页面用于展示用户过去的分类记录，采用列表布局展示每条记录的缩略图、分类标签、置信度和分类时间。历史记录数据使用鸿蒙的\texttt{@ohos.data.relationalStore}关系型数据库进行持久化存储。

历史记录页面的实现如下：

\begin{lstlisting}[language=Java, caption={历史记录页面实现}]
import relationalStore from '@ohos.data.relationalStore'

@Entry
@Component
struct HistoryPage {
  @State historyList: HistoryItem[] = []

  aboutToAppear() {
    this.loadHistory()
  }

  async loadHistory() {
    const store = await relationalStore.getRdbStore(getContext(this), {
      name: 'classifier.db', securityLevel: 1
    })
    const resultSet = await store.querySql(
      'SELECT * FROM history ORDER BY timestamp DESC')
    const items: HistoryItem[] = []
    while (resultSet.goToNextRow()) {
      items.push({
        id: resultSet.getLong(0),
        imageUri: resultSet.getString(1),
        label: resultSet.getString(2),
        confidence: resultSet.getDouble(3),
        timestamp: resultSet.getString(4)
      })
    }
    resultSet.close()
    this.historyList = items
  }

  build() {
    Column() {
      Text('分类历史')
        .fontSize(22)
        .fontWeight(FontWeight.Bold)
        .margin({ top: 20, bottom: 16 })

      List({ space: 12 }) {
        ForEach(this.historyList, (item: HistoryItem) => {
          ListItem() {
            Row() {
              Image(item.imageUri)
                .width(60)
                .height(60)
                .objectFit(ImageFit.Cover)
                .borderRadius(8)
              Column() {
                Text(item.label)
                  .fontSize(16)
                  .fontWeight(FontWeight.Medium)
                Text('置信度: ' + (item.confidence * 100).toFixed(1) + '%')
                  .fontSize(13)
                  .fontColor('#666666')
                  .margin({ top: 4 })
                Text(item.timestamp)
                  .fontSize(12)
                  .fontColor('#999999')
                  .margin({ top: 2 })
              }
              .alignItems(HorizontalAlign.Start)
              .margin({ left: 12 })
            }
            .width('100%')
            .padding(12)
            .backgroundColor('#FFFFFF')
            .borderRadius(10)
          }
        })
      }
      .width('100%')
      .layoutWeight(1)
    }
    .width('100%')
    .height('100%')
    .padding({ left: 16, right: 16 })
    .backgroundColor('#F0F0F0')
  }
}
\end{lstlisting}

历史记录的保存操作在分类完成后执行，将分类结果写入数据库：

\begin{lstlisting}[language=Java, caption={历史记录保存}]
async saveHistory(label: string, confidence: number, imageUri: string) {
  const store = await relationalStore.getRdbStore(getContext(this), {
    name: 'classifier.db', securityLevel: 1
  })
  const timestamp = new Date().toLocaleString('zh-CN')
  await store.executeSql(
    'INSERT INTO history (imageUri, label, confidence, timestamp) ' +
    'VALUES (?, ?, ?, ?)',
    [imageUri, label, confidence, timestamp]
  )
}
\end{lstlisting}

本节详细介绍了系统各核心模块的设计与实现过程。开发环境搭建部分完成了训练服务器和鸿蒙开发环境的配置；ResNet-50模型训练部分在CIFAR-10数据集上实现了92.3\%的分类准确率；模型轻量化优化部分通过INT8量化和知识蒸馏将模型大小分别压缩至25.2MB和9.8MB；模型转换与部署部分完成了MindSpore到MindSpore Lite的模型格式转换和项目集成；端侧推理引擎部分实现了基于NAPI的C++推理模块和完整的图像预处理流水线；应用界面部分基于ArkTS框架实现了主页面、图片选择、结果展示和历史记录等功能。各模块协同工作，构成了完整的端侧智能图像分类应用系统。
